\section{Выводы}
В работе реализован обмен между двумя отдельными программами через неименованные каналы. Родитель управляет вводом и последовательностью запросов, а дочерний процесс выполняет проверку простоты и запись составных чисел в файл.

Разделение команд и ответов между двумя каналами позволило согласовать завершение процессов: после статуса остановки родитель больше не передаёт команды. Начальное подтверждение готовности исключает начало обмена до успешного открытия выходного файла.

Корректность работы зависит не только от передачи данных, но и от времени жизни дескрипторов. Закрытие ненужных концов каналов обеспечивает обнаружение конца потока, а ожидание ребёнка позволяет родителю получить окончательный результат его работы. RAII уменьшает число путей выполнения, на которых ресурс мог бы остаться незакрытым.

Платформенный интерфейс отделяет задачу обработки чисел от различий Windows и Linux. При этом переносимость требует согласованных правил владения ресурсами, обработки ошибок и наследования дескрипторов, а не только одинаковых имён функций. Практическая проверка также показала необходимость согласовывать имя дочерней программы в коде с именем, сформированным системой сборки.
